\let\negmedspace\undefined
\let\negthickspace\undefined
\documentclass[journal]{IEEEtran}
\usepackage[a5paper, margin=10mm, onecolumn]{geometry}
%\usepackage{lmodern} % Ensure lmodern is loaded for pdflatex
\usepackage{tfrupee} % Include tfrupee package

\setlength{\headheight}{1cm} % Set the height of the header box
\setlength{\headsep}{0mm}     % Set the distance between the header box and the top of the text

\usepackage{gvv-book}
\usepackage{gvv}
\usepackage{cite}
\usepackage{amsmath,amssymb,amsfonts,amsthm}
\usepackage{algorithmic}
\usepackage{graphicx}
\usepackage{textcomp}
\usepackage{xcolor}
\usepackage{txfonts}
\usepackage{listings}
\usepackage{enumitem}
\usepackage{mathtools}
\usepackage{gensymb}
\usepackage{comment}
\usepackage[breaklinks=true]{hyperref}
\usepackage{tkz-euclide} 
\usepackage{listings}
% \usepackage{gvv}                                        
\def\inputGnumericTable{}                                 
\usepackage[latin1]{inputenc}                                
\usepackage{color}                                            
\usepackage{array}                                            
\usepackage{longtable}                                       
\usepackage{calc}                                             
\usepackage{multirow}                                         
\usepackage{hhline}                                           
\usepackage{ifthen}                                           
\usepackage{lscape}
\begin{document}

\bibliographystyle{IEEEtran}
\vspace{3cm}

\title{4.5.13}
\author{EE25BTECH11011 - Navya}
% \maketitle
% \newpage
% \bigskip
{\let\newpage\relax\maketitle}

\renewcommand{\thefigure}{\theenumi}
\renewcommand{\thetable}{\theenumi}
\setlength{\intextsep}{10pt} % Space between text and floats
\textbf{Question}:\\
Find the equations of the lines parallel to axes and passing through $(2, 3)$

\textbf{Solution:}  

The given point is  
\begin{align}
A = \begin{pmatrix}2 \\ 3\end{pmatrix}.
\end{align}


\textbf{Case 1: Line parallel to the $x$-axis}

 A line parallel to the $x$-axis has direction vector  
\begin{align}
\vec{d} = \begin{pmatrix}1 \\ 0\end{pmatrix}.
\end{align}

 A normal vector to this line is therefore  

\begin{align}
\vec{n} = \begin{pmatrix}0 \\ 1\end{pmatrix}.
\end{align}

Equation of a line in normal form is  
\begin{align}
\vec{n}^T(\vec{x} - A) = 0.
\end{align}

Substituting,
\begin{align}
\begin{pmatrix}0 & 1\end{pmatrix}
\left(\vec{x} - \begin{pmatrix}2 \\ 3\end{pmatrix}\right) &= 0
\end{align}
\begin{align}
\begin{pmatrix}0 & 1\end{pmatrix}\vec{x} &= 3
\end{align}
\textbf{Case 2: Line parallel to the $y$-axis}

 A line parallel to the $y$-axis has direction vector  
\begin{align}
\vec{d} = \begin{pmatrix}0 \\ 1\end{pmatrix}.
\end{align}
 A normal vector to this line is therefore  
\begin{align}
\vec{n} = \begin{pmatrix}1 \\ 0\end{pmatrix}.
\end{align}

Equation of the line:
\begin{align}
\begin{pmatrix}1 & 0\end{pmatrix}
\left(\vec{x} - \begin{pmatrix}2 \\ 3\end{pmatrix}\right) &= 0 \\[6pt]
\begin{pmatrix}1 & 0\end{pmatrix}\vec{x} &= 2
\end{align}


\textbf{Case 3: General line through $(2,3)$}

Let the normal vector be  
\begin{align}
\vec{n} = \begin{pmatrix}a \\ b\end{pmatrix}, \quad (a,b)\neq (0,0).
\end{align}

Equation:

\begin{align}
\vec{n}^T(\vec{x} - A) &= 0 \\[6pt]
\begin{pmatrix}a & b\end{pmatrix}
\left(\vec{x} - \begin{pmatrix}2 \\ 3\end{pmatrix}\right) &= 0 \\[6pt]
\begin{pmatrix}a & b\end{pmatrix}\vec{x} &= 2a+3b
\end{align}
\begin{figure}[h!]
    \centering
    \includegraphics[width=1.1\textwidth]{figs/fig6.png}
    \label{fig:division}
\end{figure}
\end{document}